\documentclass[
    a4paper, % Paper size, use either a4paper or letterpaper
    10pt, % Default font size, can also use 11pt or 12pt, although this is not recommended
    %unnumberedsections, % Comment to enable section numbering
]{LTJournalArticle}

\usepackage{subfig}
\usepackage{amsmath}
\usepackage{stfloats}
\usepackage[final]{hyperref}
\hypersetup{
    colorlinks=true,
    linkcolor=blue,
    citecolor=blue,
    urlcolor=blue
}
\usepackage[nameinlink, noabbrev]{cleveref}

%\runninghead{Sting Operation: WASP-135b in Transit} % A shortened article title to appear in the running head, leave this command empty for no running head

%----------------------------------------------------------------------------------------
%	TITLE SECTION
%----------------------------------------------------------------------------------------

\begin{document}
\epsfxsize=40mm                         % Size of crest
\thispagestyle{empty} % This removes the running head from the top corner
\begin{minipage}[b]{110mm}
        {\Huge\bf School of Physics \\and Astronomy
        \vspace*{17mm}}
\end{minipage}
\hfill
\begin{minipage}[t]{40mm}               
        \makebox[40mm]{
        
\includegraphics[width=40mm]{figures/UoEcrest.pdf}}
\end{minipage}
\par\noindent                                           % Centre Title, and name
\vspace*{2cm}
\begin{center}
        \huge\bf Data Science Techniques Applied To Cosmology
\end{center}
\vspace*{1cm}
\begin{center}
        \bf 
        \large
        Grace Tait\\
        \vspace{0.5cm}
        February 2026                         % Submission Date
\end{center}
\vspace*{5mm}

%----------------------------------------------------------------------------------------
%	ABSTRACT
%----------------------------------------------------------------------------------------


                                         
\begin{abstract}
Neutral hydrogen (HI) 21\,cm intensity mapping is a promising probe of the large-scale structure of the Universe, but is contaminated by foreground emission several orders of magnitude brighter than the cosmological signal. This work compares Principal Component Analysis (PCA) and Gaussian Process Regression (GPR) for foreground removal methods in simulated 21\,cm intensity mapping. PCA successfully recovered the HI power spectrum after removing three principal components, with residuals consistent with the true signal after applying minimum mode thresholds of $k_{\perp}$ = 0.1\,Mpc$^{-1}$ and $k_{\parallel}$ = 0.1\,Mpc$^{-1}$. GPR, using an exponential kernel for the HI signal and a spectral-index-modified RBF kernel for the foreground, yielded incomplete foreground removal due to poorly constrained foreground hyperparameters. While PCA outperforms GPR for this idealised simulation by leveraging the statistical structure in the data, GPR's ability to propagate uncertainties and incorporate physical priors may be a valuable asset for next-generation surveys with MeerKAT.
\end{abstract}


\vspace*{1cm}

\vspace*{3cm}

\vfill
\begin{center}
    {\bf Supervisor:} Dr. Trent Dupuy                % Change to suit
\end{center}
 \newpage


%----------------------------------------------------------------------------------------
%	TABLE OF CONTENTS
%----------------------------------------------------------------------------------------

\newpage
\thispagestyle{empty} % This removes the running head from the top corner
%\tableofcontents
\date{}
\newpage

%----------------------------------------------------------------------------------------
%	ARTICLE CONTENTS
%----------------------------------------------------------------------------------------
\setcounter{page}{1} % Set this page as the first numbered page
\twocolumn

% Main sections 
\section{Introduction}

% [1] https://science.nasa.gov/dark-energy/#:~:text=Astronomer%20Edwin%20Hubble%20confirmed%20that,they%20confirmed%2C%20is%20really%20expanding.
% [1] https://www.pnas.org/doi/10.1073/pnas.15.3.168
% [2] https://arxiv.org/abs/astro-ph/9712212
% [3] https://arxiv.org/abs/2203.06142v1

% [1] https://arxiv.org/abs/astro-ph/9805201
The expansion of the Universe was confirmed in 1929 through observations of the redshifts of spiral galaxies (\cite{hubbleRelationDistanceRadial1929a}). At the time, it was expected that gravitational attraction would act to decelerate this expansion, a prediction consistent with the framework of general relativity for a matter-dominated Universe. However, observations of high-redshift Type Ia supernovae revealed them to be further and fainter than expected, providing the first direct evidence that the expansion of the Universe is in fact accelerating (\cite{perlmutterDiscoverySupernovaExplosion1998a}, \cite{riessObservationalEvidenceSupernovae1998}). This discovery significantly changed our understanding of the Universe. In response, the scientific community turned to a range of independent observational probes, such as the Cosmic Microwave Background (CMB), Type Ia Supernovae, and Baryon Acoustic Oscillations, to better understand this acceleration (\cite{abdallaCosmologyIntertwinedReview2022a}). 
% also talk about any information we can get from measuring the acceleration of the Universe and why it's significant
% dark matter and dark energy
% talk about the current standard cosmological model

% [4] https://arxiv.org/html/2312.08542v1?st_source=ai_mode#:~:text=One%20of%20the%20most%20important,et%20al.%2C%202022)%20.
% [5] https://arxiv.org/abs/1807.06209
% [6] https://arxiv.org/abs/2112.04510
% [7] https://arxiv.org/abs/2106.15656
% [8] https://arxiv.org/abs/2308.10954
% [9] https://arxiv.org/abs/2107.10291
% [10] https://arxiv.org/abs/2103.01183
Since this breakthrough, as instrumental technology advanced, the volume and precision of cosmological data increased significantly and led to high precision in cosmological parameters. While the majority of observations agree with the predictions of the standard cosmological model, there are discrepancies between these methods, known as cosmological tensions (\cite{leizerovichTensionsCosmologyDiscussion2024a}). The most prominent being the discrepancy between the Hubble constant $H_0$ determined from the CMB (\cite{collaborationPlanck2018Results2020a}) not being in agreement with $H_0$ determined from Type Ia supernovae and Cepheids (\cite{riessComprehensiveMeasurementLocal2022a}). Although this discrepancy has been discussed at length in literature, there is no agreement in the scientific community on the source (\cite{freedmanMeasurementsHubbleConstant2021}, \cite{riessLocalValueH$_0$2024}, \cite{schoneberg$H_0$OlympicsFair2022}). Due to the increased precision in the current cosmological probes becoming close to reaching its maximum, it is prudent to look for new and independent cosmological probes that could either confirm or deny the discrepancies found (\cite{valentinoRealmHubbleTension2021}).

% [11] https://www.cfa.harvard.edu/research/topic/large-scale-structure?st_source=ai_mode (first sentence but havent cited it yet)
% [12] https://arxiv.org/pdf/2206.01579
% or https://academic.oup.com/mnras/article/423/4/3430/1747727?st_source=ai_mode ??
% [13] https://academic.oup.com/mnras/article/355/4/1339/993455?login=false
% [14] https://academic.oup.com/mnras/article/541/1/476/8161700
Probing the large scale structure (LSS) of the Universe is crucial in providing information about the very early Universe, the distribution of dark matter, galaxy formation, and measuring the Universe's acceleration. Traditionally, this involved measuring the redshift and angular position of galaxies individually, then creating a catalogue and 3D map in the optical or near-infrared (\cite{cunningtonHIIntensityMapping2022}). While this method led to constraints on dark energy, gravity, and the initial conditions of the Universe, it often requires the difficult detection of individual galaxies which becomes expensive. An alternative and novel probe for the LSS of the Universe is Neutral Hydrogen (HI) Intensity Mapping (\cite{battyeNeutralHydrogenSurveys2004}) which overcomes these challenges using the cumulative emission of HI flux in the sky (\cite{mazumderHiIntensityMapping2025}). Hence, it can map large areas of the sky very quickly. 

% [14] https://arxiv.org/pdf/0902.3091
% [15] https://arxiv.org/pdf/0801.1677
% [16] https://iopscience.iop.org/article/10.1086/504972
% [17] https://journals.aps.org/prd/abstract/10.1103/PhysRevD.76.083005
Following the epoch of reionization, majority of HI exists inside of galaxies due to self-shielding. This makes HI an excellent tool for probing how matter is distributed throughout the Universe. The 21\,cm line arising from the hyperfine spin-flip transition of HI is the dominant spectral line in radio astronomical observations at frequencies below 1420MHz (\cite{peterson21CmIntensity2009}). By observing this redshifted line with radio telescopes, and due to the fact it is an isolated transition, a 3D map of the LSS of the Universe can be constructed providing the radial coordinate along the line-of-sight through the frequency dependence of the redshift. This is a significant advantage over probes such as the CMB, which maps a single 2D surface (\cite{loebPossibilityPreciseMeasurement2008a}). Furthermore, this line can in principle be observed across a wide range of cosmic epochs, from the dark ages through the epoch of reionisation to the present day (\cite{shapiro21CmBackground2006}), \cite{lewis21CmAngularpower2007a}), enabling access to the large co-moving volumes at high redshift where the majority of the Universe's structure resides. HI intensity mapping therefore offers a powerful probe of the large-scale matter distribution.

% the faintness of HI emission lines make it challenging to resolve individual galaxies but for cosmology we are interested in the 
% talk about meerkat arrays. mapping out stuff with sufficient baseline/sensitivity

% [18] https://ui.adsabs.harvard.edu/abs/1999A%26A...345..380S/abstract
% [19] https://arxiv.org/pdf/astro-ph/0302099 
One of the main challenges for HI intensity mapping is the presence of foreground emission, such as synchrotron and free-free emission from the Milky Way (\cite{shaverCanReionizationEpoch1999a}), which dominates at this wavelength. The foregrounds are expected to be spectrally and spatially smooth in comparison to the rapid variations in frequency of the HI signal, which means that even though the HI signal is several orders of magnitude fainter than the foreground, the signal can be separated out (\cite{ohForegrounds21cmObservations2003a}). 

% [2] https://academic.oup.com/mnras/article/488/4/5452/5531337
% [3] https://research.manchester.ac.uk/en/publications/interference-from-global-navigation-satellites-in-future-hi-inten/
Beyond foreground contamination, HI intensity mapping faces additional challenges including thermal noise from the receiver system, frequency-dependent beam effects that cause mixing between angular and spectral modes (\cite{cunningtonImpactForegroundsIntensity2019}), and radio frequency interference from human-made transmitters (\cite{harperInterferenceGlobalNavigation2018}). Current and next-generation radio telescopes such as MeerKat and the Square Kilometre Array (SKA; \cite{groupCosmologyPhase12020}) are expected to map HI intensity across large volumes of the Universe with unprecedented sensitivity (\cite{hothiComparingForegroundRemoval2021}), making the development and validation of robust foreground removal techniques an increasingly urgent priority.
% [7] https://arxiv.org/abs/1811.02743
% [8] https://arxiv.org/abs/1501.03989

% [4] https://arxiv.org/abs/2511.13328?st_source=ai_mode#:~:text=Our%20goal%20is%20to%20evaluate,%5Bastro%2Dph.CO%5D
% [5] https://iopscience.iop.org/article/10.1086/310631
% [6] https://arxiv.org/abs/2010.15843
% [8] https://academic.oup.com/mnras/article/500/2/2264/5979792
To address these challenges, a range of foreground removal techniques have been proposed and tested in the literature, including polynomial fitting along the frequency axis (\cite{gkogkouForegroundRemovalHI2026}), Independent Component Analysis (\cite{hothiComparingForegroundRemoval2021}), Wiener filtering (\cite{tegmarkHowMakeMaps1997}), and more recently deep learning approaches (\cite{makinenDeep21DeepLearning2021}). Each method represents a trade-off between the assumptions required and the accuracy of the recovered signal.

This work applies and compares two foreground removal methods: Principal Component Analysis (PCA) and Gaussian Process Regression (GPR). PCA is a blind statistical method which exploits the spectral smoothness of the foregrounds in comparison to the rapid fluctuations of the HI signal. GPR is a method that explicitly models the signal covariance. These methods are applied to simulated 21\,cm intensity mapping data with performance being evaluated through the recovery of the HI power spectrum, in order to assess the relative strengths and limitations of each approach.

The paper is organised as follows: \Cref{sec:simulated_dataset} introduces the simulated dataset. \Cref{sec:foreground_removal_methods} describes the PCA and GPR foreground removal methodologies. \Cref{sec:results} presents the recovered HI power spectrum results for both methods. \Cref{sec:conclusion} provides a summary and conclusions.

\section{Simulated Dataset}

\begin{figure}[h]
    \centering
    \includegraphics[width=1\linewidth]{outputs/week3/hydrogen_map_wo_foreground.pdf}
    \caption{A cosmological neutral hydrogen (HI) distribution generated from theory without any foregrounds.}
    \label{fig:map_wo_foreground}
\end{figure}

\begin{figure}[h]
    \centering
    \includegraphics[width=1\linewidth]{outputs/week3/hydrogen_map_w_foreground.pdf}
    \caption{A cosmological neutral hydrogen (HI) distribution generated from theory with foregrounds.}
    \label{fig:map_w_foreground}
\end{figure}

\begin{figure}[h]
    \centering
    \includegraphics[width=1\linewidth]{outputs/week3/frequency_vs_temperature.pdf}
    \caption{A plot of frequency versus temperature. The foregrounds and HI signal are different in frequency structure. The foregrounds are 4 orders of magnitude larger than the HI signal.}
    \label{fig:frequency_vs_temperature}
\end{figure}

% [1] https://meer21cm.readthedocs.io/en/latest/
A cosmological neutral hydrogen (HI) distribution was generated from theory using the meer21cm package [1] by defining some input model. This gave a survey area and frequency range which gives the cosmological redshifts. These maps generated with just the 21cm signal and the 21cm signal with foreground can be seen in \cref{fig:map_w_foreground} and \cref{fig:map_wo_foreground}. The colour scales for both distributions are different with the foregrounds being four orders of magnitude larger than the 21cm signal. It can also be seen that they are different in frequency structure. It is important to note that this is mock data and not a real observation which would contain telescope beam smoothing, foreground contamination, and thermal noise.

A plot of the frequency versus temperature can be seen in \cref{fig:frequency_vs_temperature}. The foregrounds and HI signal are different in frequency structure. The foregrounds are 4 orders of magnitude larger than the HI signal.

% write more!! maybe
From these HI maps, the power spectrum can be extracted. This power spectrum can be flattened to 1D. See a plot which shows the power spectrum of the data compared to the input data.
\section{Foreground Removal Methods}
\label{sec:foreground_removal_methods}

This section describes the theory and the methodology of the two foreground removal methods: Principal Component Analysis (PCA) and Gaussian Process Regression (GPR). For both methods, the recovered power spectra are used to diagnose and assess the foreground cleaning and the residual signal. 

The spherically averaged 1D power spectrum takes the average of all the different $k$ modes, assuming statistical isotropy where the Universe looks the same in all directions on average. The power spectrum spans many orders of magnitude across the $k$ range, particularly where the foreground contributes significant amounts of power at low $k$, therefore the log is taken. This makes it easier to see the shape and relative difference across scales. However, directional information about the signal is lost through the averaging process. 

To recover this valuable information, the cylindrical power spectrum is also computed, providing more information on the total signal by separating modes along the line of sight ($k_{\parallel}$) from transverse modes ($k_{\perp}$). This power spectrum separates the wave-vector $k$ into its components $k_\perp$, corresponding to the spatial fluctuations across the sky which are perpendicular to the line-of-sight, and $k_\parallel$, corresponding the the frequency-dependent fluctuations along the line-of-sight. 

The 1D spherically averaged power spectrum was computed by binning modes into shells of width $\Delta k$ across the range 0 - 1.2\,Mpc$^{-1}$. The cylindrical power spectrum was computed by bining modes into a grid of $k_{\perp}$ $\epsilon$ $[0.01, 0.5]$\,Mpc$^{-1}$ and $k_{\parallel}$ $\epsilon$ $[0.01, 1.0]$\,Mpc$^{-1}$, with 20 and 10 bins respectively. Both were computed using the meer21cm package. Together, these two diagnostics provide a complete picture of foreground removal performance across all scales.

\subsection{Principal Component Analysis}

\begin{figure*}[h!]
    \centering
    \includegraphics[width=0.7\linewidth]{outputs/week4/eigenvectors.pdf}
    \caption{A plot showing the eigenvector modes versus frequency of the mock data with the total signal consisting of the 21\,cm signal and foreground.}
    \label{fig:eigenvectors}
\end{figure*}

\begin{figure*}[h!]
    \centering
    \includegraphics[width=0.65\linewidth]{outputs/week4/eigenvalues.pdf}
    \caption{A plot showing the eigenvalues versus eigenvalue index of the mock data with the total signal consisting of the 21\,cm signal and foreground.}
    \label{fig:eigenvalues}
\end{figure*}

\begin{figure*}[t!]
    \centering
    \includegraphics[width=1\linewidth]{outputs/week9/1d_log_power_spectrum_[0, 0.03]_None_both.pdf}
    \caption{The mean cylindrical power spectra of 100 independent realisation from for the HI signal, the PCA residual total signal (HI signal and foreground) after foreground removal, and the ratio between the PCA residual signal and the HI signal, from left to right.}
    \label{fig:cylindrical_power_spectrum_None_None}
\end{figure*}

% [1] https://arxiv.org/pdf/2206.01579
% [2] https://link.springer.com/book/10.1007/b98835 (last sentence but not cited yet)
Principal Component Analysis (PCA) is a blind foreground removal method which requires no prior knowledge of the foreground. This technique takes advantage of the foreground being several orders of magnitude larger than the 21\,cm signal and having a highly correlated frequency structure (\cite{cunningtonHIIntensityMapping2022}).

The data can be described as a matrix $d$ which is a linear system with dimensions $N_\theta \times N_\nu$:

\begin{equation}
    \label{eq:data}
    d = As + d_{clean},
\end{equation}

where $A$ is the mixing matrix and $s$ are the $N_{fg}$ separable foreground components identified by projecting the mixing matrix along the data ${s = A^Td}$. The mixing matrix is extracted from the eigen-decomposition of the frequency covariance matrix of the mean-centred data. This frequency-frequency matrix is defined by 

\begin{equation}
    C = \frac{(wd)^T(wd)}{N_\theta-1},
\end{equation}

where $w$ are the inverse variance weights recast into $N_\theta \times N_\nu$ matrices. The eigen-decomposition of the covariance matrix is given as

\begin{equation}
    C = V\Lambda V^T,
\end{equation}

where $V$ are the eigenvectors ordered by decreasing eigenvalue, and $\Lambda$ is the diagonal matrix of eigenvalues. The first eigenvectors, associated with the largest eigenvalues, correspond to the modes of greatest variance and are dominated by the spectrally smooth foreground emission $N_{fg}$. The first $N_{fg}$ eigenvectors supply the columns of the mixing matrix $A$, since the foreground dominates the largest variance modes. The foreground estimate $As$ can therefore be written explicitly in terms of the eigenvectors as

\begin{equation}
    As = \sum_{n=1}^{N_{fg}}v_nv_n^Td
\end{equation}

where $v_n$ are the foreground eigenvectors being projected out. By substituting this into \Cref{eq:data}, rearranging for $d_{clean}$ gives the foreground-cleaned data: 

\begin{equation}
    d_{clean} = d - As = (I - \sum_{n=1}^{N_{fg}}v_nv_n^T)d,
\end{equation}

where $I$ is the identity matrix. The underlying assumption is that the foreground emission dominates the largest variance modes, so projecting them out leaves a residual signal dominated by the 21\,cm contribution. 

%While effective, this foreground cleaning can also remove some HI signal modes which overlap with the foreground, or leave behind some residual foreground contamination in the cleaned data.

A key strength of PCA is that is requires no prior knowledge of the foreground and it is entirely data driven. It is also computationally efficient, requiring only a single eigendecomposition of the covariance matrix.

However, PCA has two inherent limitations. First, it is a blind method, which means it cannot distinguish between foreground and HI signal modes that share similar spectral structure. At low $k_{\parallel}$, where the HI signal has long frequency correlation scales comparable to the foreground, PCA inevitably removes some genuine cosmological signal alongside the foreground modes. Second, the choice of $N_{FG}$ is not uniquely determined by the data and requires a judgement based on the eigenvalue spectrum. An insufficient number of components leaves residual foreground, while too many risk subtracting genuine HI signal. More rigorous approaches such as cross-validation or Bayesian model selection could determine $N_{FG}$ more objectively. 

\subsubsection{Foreground Removal}

A HI intensity map was generated for a total signal consisting of the 21\,cm signal and the foregrounds. In order to prevent the results being subject to cosmic variance and to quantify the statistical uncertainty on the recovered power spectrum, the foreground cleaning procedure was repeated across 100 independent realisations of the mock data, generated by varying the random seeds used in the simulation. Each realisation produces a different but statistically equivalent realisation of the HI signal, and the spread across realisations reflects the sample variance inherent to a single survey volume. The mean and standard deviation of the data across all realisations were used for plotting the power spectra. This approach provides a more robust assessment of PCA performance than a single realisation alone and allows the consistency of the foreground removal method to be evaluated across different signal realisations.

The eigenvectors and eigenvalues of the total signal covariance are shown in \Cref{fig:eigenvectors} and \Cref{fig:eigenvalues}. The first two eigenvectors are visibly spectrally smooth, consistent with the expected frequency structure of synchrotron and free-free foreground emission. The remaining eigenvectors exhibit rapid frequency oscillations characteristic of the 21\,cm signal and noise. The eigenvalue spectrum shows a sharp drop of approximately 10-11 orders of magnitude after the third eigenvalue, after which the spectrum plateaus at a much lower level. This separation provides justification for the choice of $N_{FG} = 3$, as the first three modes capture the foreground-dominated variance while the remaining modes are consistent with the HI signal. This break in the eigenvalue spectrum is the key diagnostic used to select $N_{FG}$.

%The first few eigenvectors are spectrally smooth, exhibiting the smooth characteristic expected from the foreground. Additionally, the first three eigenvalue modes are the largest and correspond to the foreground. Therefore, $N_{FG} = 3$ principal components were removed from the total signal to perform the foreground cleaning. This value was chosen to maximise foreground suppression while minimising the loss of the HI signal. 

%To recover the desired cosmological information, the equation is rearrange to isolate the residuals: 

%\begin{equation}
%    \epsilon = \textbf{d} - \textbf{As}.
%\end{equation}

\subsubsection{Cylindrical Power Spectrum}

To evaluate the foreground removal, the true HI power spectrum, the residual signal, and their ratio were plotted, seen in \Cref{fig:cylindrical_power_spectrum_None_None}. By inspecting the ratio panel, the components of the wave-vector experiencing signal loss were identified. For the majority of the panel the pixels are red, indicating a ratio of unity, which means that the HI power spectrum has successfully been recovered. However, at low $k_{\parallel}$ and $k_{\perp}$, this ratio goes below unity, revealing signal loss. 

By inspection of where the blue pixels were, an initial minimum threshold of $k_{\parallel} = 0.025$\,Mpc$^{-1}$ and $k_{\perp} = 0.1$\,Mpc$^{-1}$ was tested, followed by different paris of values. The aim was to find the pair of values for $k_{\parallel}$ and $k_{\perp}$ which balanced removing enough foreground without causing too much HI signal loss. Once the optimal threshold was found, foreground cleaning was performed again using the optimal minimum thresholds of wavenumber components, the 1D power spectrum was extracted again and compared to theory.


\subsection{Gaussian Process Regression}

% [1] https://link.springer.com/chapter/10.1007/978-3-540-28650-9_4
% [2] https://mc-stan.org/docs/2_38/stan-users-guide/gaussian-processes.html#:~:text=Unlike%20a%20simple%20multivariate%20normal,mean%20function%20and%20covariance%20function.
In Gaussian Process Regression (GPR), the various components of the signal, which are the foreground and 21\,cm signal, are modelled as realisations of a Gaussian Process (GP; \cite{rasmussenGaussianProcessesMachine2008}). A GP is a Gaussian distribution extended to infinite dimensions. While a standard multivariate Gaussian distribution describes a finite set of variables, a GP expands this to an infinite number of variables. In this context, instead of describing discrete points, the distribution describes continuous functions over a range of frequencies $\nu$. Consequently, a GP provides a probability distribution over a space of possible functions. 

In contrast to multivariate Gaussian distributions, which are defined by a mean vector and covariance matrix, a GP is defined by a mean function and a covariance function. The covariance function is defined by a kernel, which determines the relationship between different points in the data and specifies which functions are likely under the GP prior. This determines the generalisation properties of the model.

%In contrast to multivariate Gaussian distributions, which is defined by a mean vector and covariance matrix, a GP is defined by a mean function and a covariance function. In this case, each element of the mean is a function of frequency and the covariance function is defined by a kernel which means they are easy to sample from. A kernel, also known as a covariance function or kernel function, defines the shape of the functions by determining the relationship between different points in the data and is the set of functions considered for the data. It specifies which functions are likely under the GP prior. This, then, determines the generalisation properties of the model.

GPR operates in a similar way to regression with parametric models, with the main distinction being that a prior is put on the the kernel rather than on the specific model coefficients, where conventionally the prior has a mean of zero. GPR is a non-parametric regression technique, meaning it is able to estimate relationships between variables without assuming a predefined function form. Instead, the model is constructed directly from the data, making it useful for complex, unknown, or non-linear datasets. 

% [3] https://arxiv.org/pdf/2102.01697
Another advantage of GPs is that it is stochastic, which mean that each realisation is random and can take any functional form consistent with the mean and kernel function. Assuming only Gaussian noise is present, they also have a well defined marginal likelihood function, so it is easy to compute the probability of data conditioned on a given value of the mean and kernel functions (\cite{lugerMappingStellarSurfaces2021}). %This can be maximised to find the optimal values of the model parameters or be used to evaluate the posterior which is very useful for inference problems.

However, GPs require a larger amount of data to achieve the same level of uncertainty as a parametric model, pose a higher risk of overfitting, are slower to compute, and can be difficult to interpret.

The assumption is made that the data, and each of its components, are drawn from a GP with a kernel function and zero mean, which is consistent with the data being mean-centred. Since the total data is composed of separate GP components, the total kernel can be decomposed as

\begin{equation}
    K_{tot} = K_{fg}+K_{HI},
\end{equation}

where $K_{fg}$ and $K_{HI}$ are the foreground and HI kernels respectively.

Given this decomposition, the GP conditional mean provides the optimal estimate of the foreground signal given the total observed data. The foreground contribution is estimated as

\begin{equation}
    \hat{f} = K_{fg}(K_{fg}+K_{HI})^{-1}d.
\end{equation}

The cleaned residual map is then obtained by subtracting this estimate from the total data:

\begin{equation}
    d_{clean} = d - \hat{f} = [I - K_{fg}(K_{fg}+K_{HI})^{-1}]d
\end{equation}

\begin{equation}
    d_{clean} = K_{HI}(K_{fg}+K_{HI})^{-1}d
\end{equation}

The cleaned map retains only the component of the data best explained by the HI kernel, with the foreground contribution projected out. The quality of this cleaning therefore depends critically on how well the fitted kernels describe the true covariance structures of each component.

%Since the data is made up of separate components which are all GPs, the kernel function can also be split into separate components: the foreground kernel and the 21\,cm signal kernel. Since the foreground model is a GP, it has an expectation value and covariance. The expectation value is a prediction of the foreground signal and is dependent on both the foreground kernel, the total kernel of the data, and the data. While, the covariance provides information on the uncertainty on this prediction. Therefore, it is important there are suitable kernels for all components of the data. 

%By making an educated estimate of the true foreground signal with an appropriate kernel, the foreground signal can be modelled. This can be removed from the data, by subtracting the expectation value of the foreground from the data, and the residual signal, ideally left with the instrumental noise the the desired 21\,cm signal, is achieved. 

\subsubsection{Kernel functions}

% [1] https://www.cs.toronto.edu/~duvenaud/cookbook/
% [2] https://www.repository.cam.ac.uk/items/a54f0711-f777-4ee1-a685-3b76dae5ad91
% [3] https://gaussianprocess.org/gpml/chapters/RW.pdf
The selection of kernel functions is critical as it significantly affects the generalisation properties of the GP model. The choice of kernel functions is data driven, as they must be chosen to best describe the covariance of the data, along with the best-fitting hyperparameters, which are the parameters of the kernel functions. One commonly used function is the Matern kernel (\cite{rasmussenGaussianProcessesMachine2008}):

\begin{equation}
    K(\nu, \nu') = \sigma^2 \frac{2^{1-\eta}}{\Gamma(\eta)}(\sqrt{2\eta}\frac{|\nu - \nu'|}{l})^\eta K_\eta (\sqrt{2\eta}\frac{|\nu - \nu'|}{l}),
\end{equation}

where $\Gamma$ is the gamma function, $K_\eta$ is the modified Bessel function. This kernel has three hyperparameters which be need to optimised:

\begin{itemize}
    \item \textbf{Variance $\sigma^2$:} this describes the amplitude or "power" of the signal. Since the foregrounds are significantly brighter than the HI signal, $\sigma^2$ is expected to be much larger for the foreground component (see \Cref{fig:frequency_vs_temperature}).
    \item \textbf{Lengthscale $l$:} this describes the scale of frequency correlations. A larger $l$ indicates the data points are more correlated over wide frequency ranges. A larger lengthscale is expected for foregrounds since they are much more correlated in frequency (see \Cref{fig:tot_fg_hi_noise_covariance}).
    \item \textbf{Spectral parameter $\eta$:} this determines the smoothness of the function. Since the foregrounds are smooth in frequency and the HI signal is more Gaussian, $\eta$ is expected to be larger for the foregrounds than for the HI signal (see \Cref{fig:frequency_vs_temperature}).
\end{itemize}

By choosing different values for the spectral parameter $\eta$, the Mat\'ern kernel can simplify. This project will use two choices for this hyperparameter:

% [4] https://arxiv.org/abs/2010.15892
% [5] https://academic.oup.com/mnras/article/493/2/1662/5727337#sec3
% [6] https://arxiv.org/pdf/2105.12665
% [7] https://academic.oup.com/mnras/article/478/3/3640/4995230
\begin{itemize}

    \item \textbf{$\eta \rightarrow \infty$:} the Mat\'ern kernel simplifies to the Radial Basis Function (RBF), also known as the squared exponential or Gaussian kernel. This represents an extremely smooth signal which makes it the ideal choice for modelling the foreground component. \cite{kernGaussianProcessForeground2020}, \cite{mertensImprovedUpperLimits2020} and \cite{soaresGaussianProcessRegression2022} use this kernel to model their smooth foreground component. This has the form:

    \begin{equation}
        \label{eq:RBF}
        K(\nu, \nu') = \sigma^2 exp(\frac{(\nu - \nu')^2}{2l^2}).
    \end{equation}

     %\item \textbf{$\eta = 5/2$:} the Matern kernel is not quite as smooth in frequnecy as the RBF kernel but is still relatively smooth. This kernel was used to model the smooth foreground component by [7].
    
    %\item \textbf{$\eta = 3/2$:} the Matern kernel is even less smooth in frequency and can be used to describe foreground components with moderate frequency smoothness. This was done by \cite{mertensImprovedUpperLimits2020} and [7].
    
    \item \textbf{$\eta = 1/2$:} the Mat\'ern kernel is the least smooth in frequency and simplifies to the exponential function. Since it is not smooth at all, it is a strong match for signals that vary significantly with frequency, such as the HI signal. This is used to describe the HI signal by \cite{kernGaussianProcessForeground2020}, \cite{mertensImprovedUpperLimits2020}, and \cite{soaresGaussianProcessRegression2022} and has the form: 

    \begin{equation}
        \label{eq:exp}
        K(\nu, \nu') = \sigma^2exp(\frac{|\nu - \nu'|}{l}).
    \end{equation}
    
\end{itemize}

The versatility of the Mat\'ern kernel makes it a strong tool for various signals. By adjusting its hyperparameters, each component of the data can be effectively modelled.

\begin{figure}[t!]
    \centering
    \includegraphics[width=1\linewidth]{outputs/week6/exp_kernel_lengthscales_2.pdf}
    \caption{Frequency covariance of the exponential kernel with a low variance and low lengthscale ($\sigma = 1$, $l = 1$) seen on the left, and another with a high variance and high lengthscale ($\sigma = 100$, $l = 100$) seen on the right. This exponential kernel is used to model the 21\,cm cosmological signal.}
    \label{fig:exp_kernel_lengthscales}
\end{figure}

\begin{figure}[t!]
    \centering
    \includegraphics[width=1\linewidth]{outputs/week6/rbf_kernels_lengthscales_2.pdf}
    \caption{Frequency covariance of the radial-basis function (RBF) kernel with a low variance and low lengthscale ($\sigma = 1$, $l = 1$) seen on the left, and another with a high variance and high lengthscale ($\sigma = 100$, $l = 100$) seen on the right.}
    \label{fig:rbf_kernels_lengthscales}
\end{figure}

\begin{figure}[t!]
    \centering
    \includegraphics[width=1\linewidth]{outputs/week6/rbf_kernel_w_sp_indx_vs_foreground_covariance.pdf}
    \caption{Frequency covariance of the spectral index modified radial-basis function (RBF) kernel with high lengthscale and a spectral index ($l = 400$, $\gamma=-2.7$) seen on the left, and the true foreground covariance seen on the right. This spectral index modified RBF kernel is used to model the foreground signal.}
    \label{fig:rbf_kernel_w_sp_indx_vs_foreground_covariance}
\end{figure}

To demonstrate the behaviour of the exponential and RBF kernels, both are presented in two cases: one with a low variance and low lengthscale ($\sigma = 1$, $l = 1$), and another with a high variance and high lengthscale ($\sigma = 100$, $l = 100$), shown in \Cref{fig:exp_kernel_lengthscales} and \Cref{fig:rbf_kernels_lengthscales}. Increasing $l$ broadens the diagonal stripe, raising the correlation of the signal. Increasing $\sigma$ raises the amplitude of the covariance. By comparing these plots to the true covariance structures (\Cref{fig:tot_fg_hi_noise_covariance}), it is clear that the exponential kernel with small $\sigma$ and $l$ best describes the HI signal, while the RBF kernel with large $\sigma$ and $l$ best describes the foreground.

However, the standard RBF kernel has diagonal symmetry not seen in the true foreground covariance. To better reproduce the foreground's frequency-dependent amplitude scaling, which is physically motivated by the power-law spectrum of the synchrotron emission, a spectral index parameter $\gamma$ is introduced such that 

%We can see that the RBF function has a much wider diagonal in both cases. A wider diagonal means that the data is much more correlated in frequency. A smaller diagonal means that the data points are less correlated in frequency and are mainly correlated with themselves and the points directly around them. 

%By inspection, we can tell that the exponential kernel function with low $\sigma$ and low $l$ would be the best fit for the 21\,cm signal while the RBF function with a high $\sigma$ and high $l$ would best fit the foreground component. However, while the amplitude of the RBF covariance matches that of the foreground, the RBF covariance has diagonal symmetry that is not seen in the foreground. In order for the shape of the RBF function to match this foreground signal better, we can add a spectral index parameter $\gamma$ so the signal can scale according to frequency 

\begin{equation*}
    T(f) = T_0(f_0) * (\frac{f}{f_0}) ^ \gamma.
\end{equation*}

This modifies the RBF kernel to

\begin{equation}
    \label{eq:RBF_sp_indx}
    K(\nu, \nu') = \sigma^2 exp(\frac{(\nu - \nu')^2}{2l^2})(\frac{\nu}{\nu_0})^\gamma(\frac{\nu'}{\nu'_0})^\gamma,
\end{equation}

where $\nu_0$ and $\nu'_0$ are the first frequency values in the range. While previous studies including \cite{kernGaussianProcessForeground2020}, \cite{mertensImprovedUpperLimits2020}, and \cite{soaresGaussianProcessRegression2022} use the standard RBF kernel for foreground modelling, this work uses the spectral-index-modified form of \Cref{eq:RBF_sp_indx}, allowing the foreground amplitude to vary with frequency in a physically motivated way. The application of this modification can be seen in \Cref{fig:rbf_kernel_w_sp_indx_vs_foreground_covariance}, where the spectral-index-modified RBF kernel better reproduces the gradient structure of the true foreground covariance.

%The application of this can be seen in \Cref{fig:rbf_kernel_w_sp_indx_vs_foreground_covariance} where we can see that this addition for the RBF kernel to mimic the reality that the foregrounds are much brighter at lower frequencies. 

% [1] https://arxiv.org/pdf/2105.12665
%The goal is to find the kernel functions which best mimic the different components of the data which are the foreground and 21\,cm signal. 

%There are a couple of ways to handle noise in GPR. Since the noise is so small in amplitude and so highly correlated in frequency, it can be treated as indistinguishable noise as done by \cite{soaresGaussianProcessRegression2022}. This assumes there is either no noise present in the data, or the noise is indistinguishable from the other signals which means that the covariance functions of the other signals describes and captures the noise. It is important to noise that this data is idealised and real data would likely results in much more complex instrumental noise which would require a more rigorous treatment.

\subsubsection{Hyperparameter optimisation and model selection}

Model selection for GPR is a two step process. The first step is to choose the kernel function that best models the data, and the second is to optimise the hyperparameters. This is done with a Bayesian approach by computing the marginal likelihood and selecting the model that maximises it. 

%This represents the probability of the data given some hyperparameters $\theta$. By integrating the likelihood $\mathcal{L}$ multiplied by the prior over the GP realisations, the Log-Marginal Likelihood (LML) can be obtained: 

%It is important to pick the kernel function which performs the best along with their most optimal hyperparameters that best model the signals. These two steps are mutually dependent. So, although we have found kernel function which match our data by inspection, we much compare the overall performance of each model. This will be done by first optimising the hyperparameters of each kernels by maximising their marginal likelihoods. This will then be compared to other optimised kernel functions. 

%To evaluate each model, the marginal likelihood, also known as the Bayesian evidence, was computed. This represents the probability of the data given some hyperparameters $\theta$. By integrating the likelihood $\mathcal{L}$ multiplied by the prior over the GP realisations $x$, the Log-Marginal Likelihood (LML) can be obtained: 

%For a kernel function with hyperparameters $\theta$, the marginal likelihood of the data can be calculated by integrating the data likelihood $\mathcal{L}$ multiplied by the prior. 

%\begin{equation}
%    p(d|\nu, \theta) = \int \mathcal{L}(d|x, \nu, \theta) p(x|\nu, \theta) dx
%\end{equation}

\begin{equation}
    \log p(d|\nu, \theta) = -\frac{1}{2} tr(CK^{-1}) - \frac{1}{2} \log|K|
\end{equation}

where $C$ is the observed frequency-frequency covariance matrix of the data, $K=K_{HI}+K_{fg}$ is the total kernel function, and $\theta$ is the set of hyperparameters. This covariance-based form is used rather than standard data-vector form because the foreground removal is performed in covariance space, which fits the kernel model to the observed covariance structure of the total signal rather than to individual data realisations. The LML can be computed analytically since each component of the data is assumed to be a GP.

To compare different kernel models, the log Bayes factor $\mathcal{Z}$ is computed as the difference in LML between two models: 

\begin{equation}
    \log \mathcal{Z} = LML_{method_1} - LML_{method_2}
\end{equation}

where a larger value of $\mathcal{Z}$ or $\log \mathcal{Z}$ means that the first model fits the data better than the second model.

% [1] https://academic.oup.com/mnras/article/493/2/1662/5727337#199912059
% [2] https://arxiv.org/pdf/2105.12665
In practice, the optimal hyperparameters can be found using standard optimisation, Monte Carlo Markov Chain (MCMC) methods as used by \cite{mertensImprovedUpperLimits2020}, or Nested Sampling as used by \cite{soaresGaussianProcessRegression2022}. Standard optimisation is computationally inexpensive but is prone to converging in local minima and provides no uncertainty estimates. Nested Sampling is designed to compute the marginal likelihood directly and handles multimodal posteriors well, but is highly computationally expensive. MCMC provides a full posterior distribution over parameters, is better at exploring the global parameter space, and is computationally intermediate between the two methods.

%On the other hand, Nested Sampling is designed specifically to calculate the marginal likelihood and is good at handling data with several separate peaks with high likelihoods. This, however, is highly complex and much more computationally expensive. MCMC Sampling provide a full posterior showing the probability distribution and correlation between parameters. They are also better at finding the global peak as it explores parameter space. This method takes much longer than standard optimisation techniques but runs quicker than Nested Sampling, thus, being a middle ground between the two methods. 

% [1] https://ui.adsabs.harvard.edu/abs/2013PASP..125..306F/abstract
In this work, the SciPy minimize function (\cite{virtanenSciPy10Fundamental2020}) was used for initial hyperparameter optimisation to efficiently compare LMLs across different combinations of algorithms, hyperparameter initial guesses, and bounds. The LMLs obtained from this initial optimisation were compared across kernel combinations to identify the best-fitting model, which was then used for MCMC sampling using the \texttt{emcee} package (\cite{foreman-mackeyEmceeMCMCHammer2013a}).

Initial guesses for the optimisation were chosen by visual comparison of the kernel functions to the true covariance structures. For example, the 21\,cm signal is tightly correlated along the diagonal (see \Cref{fig:tot_fg_hi_noise_covariance}), suggesting small $l$ and small $\sigma$ with the exponential kernel. The foreground is smooth and high-amplitude, suggesting large $l$, large $\sigma$, and a negative spectral index with the RBF kernel. Bounds were set to exclude unphysical values. For example, the spectral index was constrained to be negative, consistent with the observed decrease in synchrotron brightness with increasing frequency.

%Initial guesses must be provided for the algorithm to start on. These intial guesses are important as the algorithm is sensitive to this and an incorrect starting point may result in a final value which is unphysical. These initial guesses were taken by fitting different values of the lengthscale $l$ and variance $\sigma$ to the kernels to try to mimic the true covariance of the signal components. For example, we know the 21\,cm signal is very tightly and strongly correlated in frequency (see \Cref{fig:tot_fg_hi_noise_covariance}) so a smaller $l$ and smaller $\sigma$ would be suitable with the exponential kernel (see \Cref{fig:exp_kernel_lengthscales}). On the other hand, the foreground is very smooth in frequency and high in amplitude (see \Cref{fig:tot_fg_hi_noise_covariance}) so the RBF kernel with high $l$, high $\sigma$ and a spectral index would be most suitable (see \Cref{fig:rbf_kernel_w_sp_indx_vs_foreground_covariance}).

%In addition to this, bounds must be set for each of the hyperparameters for both of these methods. This means that the algorithms will not explore any values that are unphysical. Again, these are chosen of what would be appropriate by inspection of the kernel functions and with knowledge of the data. For example, we know that the spectral index must be negative otherwise the amplitude would increase with frequency and that is the opposite of what we see. A reasonable range around the initial guesses would be a suitable choice. 

% [2] https://joss.theoj.org/papers/10.21105/joss.00024
The 100 walkers were initialised in a tight Gaussian ball centred on the best-fit hyperparameters from the initial optimisation, with a perturbation of 0.1\% applied to each hyperparameter. A prior check was applied to ensure all walkers began within the defined flat (uniform) prior bounds. Each walker evolved by proposing moves informed by the positions of the other walkers, preferentially sampling regions of high posterior probability while retaining the ability to explore lower-probability regions, avoiding convergence to local maxima. The chains were run for 10,000 steps with the first 1,000 discarded as burn-in. The thinning factor was chosen to be approximately half the autocorrelation time of the chain, ensuring that successive samples are approximately independent. Hence, the chain was thinned by retaining every 40th steps, yielding an approximately independent set of posterior samples. Convergence of the walkers were assessed through the mean acceptance fraction, visual inspection of parameter trace plots, and the corner plots (\cite{foreman-mackeyCornerpyScatterplotMatrices2016a}).

To quantify the uncertainty on the recovered power spectrum arising from the kernel hyperparameters, 1,000 parameter sets were drawn at random from the MCMC posterior chain. From each draw, the corresponding HI and foreground kernels were constructed and foreground cleaning was applied, yielding 1,000 independent realisations of the cleaned map. The 16th, 50th, and 84th percentiles were then computed across these realisations at each $k$, corresponding to the median recovered power spectrum and its 1$\sigma$ interval. These represent the uncertainties on the residual power spectrum, representing how sensitive the foreground removal is to the uncertainty in the fitted kernel parameters. 

The MCMC sampling was performed with a single fixed random seed, meaning that while hyperparameter uncertainties were obtained, cosmic variance was not accounted for. Ideally, the sampling would be repeated across 100 independent realisations of the HI signal, as done with PCA, with uncertainties derived from the resulting spread. However, due to computational and time constraints, this single-seed approach was deemed sufficient for the purposes of this work.
%\section{Methodology}

\subsection{Principal Component Analysis}

% https://www.aanda.org/articles/aa/full_html/2025/11/aa53461-24/aa53461-24.html

\begin{figure}[h!]
    \centering
    \includegraphics[width=1\linewidth]{outputs/week3/frequency_frequency_covariance.pdf}
    \caption{Frequency-frequency covariance matrix which is dominated by the foregrounds.}
    \label{fig:freq_covariance}
\end{figure}

\begin{figure}[h!]
    \centering
    \includegraphics[width=1\linewidth]{outputs//week3/HI_covariance.pdf}
    \caption{HI covariance signal which is plotted separately from the foregrounds.}
    \label{fig:HI_covariance}
\end{figure}


By generating mock data using meer21cm, the frequency-frequency covariance can be plotted for the total signal of the 21cm signal and the foregrounds, seen in \cref{fig:freq_covariance}. We can see this is dominated by the foregrounds. The HI signal plotted separately can be seen in \cref{fig:HI_covariance}.

% provide more information on what kperp and kpara represents
From this HI map, the power spectrum was extracted. This power spectrum is a function of $k_\perp$ and $k_\parallel$ and provides information on how matter is distributed in the Universe. This can be seen in \cref{fig:power_spectrum_from_hydrogen_map}. The mock data can be generated using a fixed seed for random number generation, producing a specific outcome of a stochastic model known as a realisation. By running multiple realisations, the many different but equally likely outcomes of the mock data can be generated and used to quantify uncertainty. This was done by passing in various seed numbers to obtain uncertainties on the mock data. By averaging the power spectra into a 1D spectrum, the mock data can be compared to the input theory data. 

\subsection{Gaussian Process Regression}
\section{Results}
\label{sec:results}

Here, we will present the main results which is the performance of both PCA and GPR as foreground removal methods. First, the results for PCA will be presented, secondly, GPR.

\subsection{Principal Component Analysis}

The performance of the PCA foreground removal was evaluated by comparing the recovered residual signal to the true cosmological signal through the cylindrical and spherically averaged 1D power spectra. First, we will discuss the choice of foreground principal components. Then, the results from the foreground removal which is the 1D and cylindrical power spectrum will be discussed. Finally, we will finish with a review on the overall performance of PCA.

\subsubsection{Foreground Principal Components}

The choice of $N_{FG} = 3$ was decided on based by inspection of the eigenvectors and eigenvalues. This choice of $N_{FG}$ components has a trade-off where too few components leave residual foreground, while too many mean the HI signal might be subtracted. The first three modes captured the foreground-dominated variance while the remaining modes were consistent with the cosmological signal. 

% [1] https://arxiv.org/pdf/2105.12665
The choice of $N_{FG} = 3$ is much lower than what was used in [1] which used $N_{FG} = 6$ and $N_{FG} = 7$. This is likely due to their more complex data and our data only consisting of the cosmological signal and the foreground. The results suggests that three components is a reasonable choice for this simulation, though a more rigorous approach could be taken which involves cross-validating the choice of $N_{FG}$ against the recovered power spectrum as a function of component number. 

\subsubsection{Cylindrical Power Spectrum}

The cylindrical power spectrum of the true HI power spectrum, the residual signal after PCA cleaning, and their ratio (see \Cref{fig:cylindrical_power_spectrum_None_None}) was used as an assessment of the foreground cleaning. The ratio panel provided the clearest diagnostic of the cleaning performance. At high $k_{\parallel}$, the ratio is close to unity and uniformly red, indicating that the recovered signal is in good agreement with the true HI power spectrum in these modes. This is physically expected as high $k_{\parallel}$ modes correspond to small-scale structure along the line of sight that varies rapidly with frequency and is therefore largely orthogonal to the spectrally smooth foreground components removed by PCA. 

% [1] https://academic.oup.com/mnras/article/446/4/3545/2891891
% [2] https://iopscience.iop.org/article/10.1086/668105/pdf#:~:text=This%20work%20was%20originally%20developed%20for%20PCA,underlying%20methodology%20is%20specific%20to%20astronomical%20spectra.
% [3] http://academic.oup.com/mnras/article/478/3/3640/4995230#118057951
% [4] https://arxiv.org/pdf/2407.21626
However, at low $k_{\parallel}$, the ratio drops below unity, indicating significant signal loss. Here, the 21\,cm signal shares spectral structure with the foregrounds and cannot be cleanly separated by PCA. This inevitably removes some genuine HI signal alongside the foreground components, a known limitation of blind component separation techniques [1] [2]. This signal loss at low $k_{\parallel}$ is consistent with findings in [3] and [4], who report the same characteristic suppression of modes at low $k_{\parallel}$ from PCA cleaning. 

Through inspection of the ratio panel, a minimum threshold of $k_\parallel = 0.4$\,Mpc$^{-1}$ was defined and the foreground cleaning of the signal was redone excluding modes below these values. Minimum thresholds with values below this were found to still show cosmological signal loss. While minimum threshold with values larger than this were found to as well closely match the true HI signal, the goal was to find the smallest threshold which maximised foreground cleaning and minimised cosmological signal loss. 

\subsubsection{1D Power Spectrum}

\begin{figure}[t!]
    \centering
    \includegraphics[width=1\linewidth]{outputs/week9/1d_log_power_spectrum_[0, 0.04]_None_both.pdf}
    \caption{The logarithmic spherically averaged 1D power spectrum from of the HI signal with no foreground removal (black), and the residual signal after PCA foreground removal without any minimum thresholds applied (blue) and the residual signal after PCA foreground removal with minimum thresholds of $k_\parallel = 0.4$\,Mpc$^{-1}$ (red). This was run for 100 realisations with the mean used as the datapoints and the standard deviation used as the error bars.}
    \label{fig:1d_log_power_spectrum_0_0.04_None_both}
\end{figure}

The spherically averaged 1D power spectrum with the true HI signal and the residual signal before and after applying the minimum $k$ thresholds can be seen in \Cref{fig:1d_log_power_spectrum_0_0.04_None_both}. Since foreground cleaning was run with 100 independent realisations from 100 fixed seeds, this allowed for uncertainties to be quantified for both the residual signal and the cosmological HI signal, taking into account not only the foreground cleaning but also cosmic variance. Thus, error bars are also seen on the HI signal. 



Prior to applying the minimum thresholds, the residual signal (shown by the blue line) follows the true HI signal (shown by the black line) at large $k$ but shows a clear deviation at low $k$, consistent with the signal loss identified in the cylindrical power spectrum ratio panel due to some of the genuine cosmological signal being removed with the foreground components.

After applying the minimum thresholds of and $k_\parallel = 0.4$\,Mpc$^{-1}$, the residual signal (shown by the red line) shows a stronger agreement with the 21\,cm signal across the whole $k$ range, with the curve lying within 1$\sigma$ of the HI signal error bars across the whole $k$ range. We can see that the error bars are uniform in size at large $k$ but towards lower $k$ the error bars for the residual signals grow larger, especially for the residual signal with the minimum thresholds. This is expected and represents the increased variance in the residual signal as a result from the cleaning. 

It is also important to note that the residual signal with thresholds does not extend all the way to $k = 0$\,Mpc$^{-1}$ and this is a consequence of excluding these minimum thresholds as these datapoints are completely excluded from the analysis. So, while this signal does recover the true HI signal well, there is genuine cosmological signal loss. 

\subsubsection{Overall Performance}

Overall, the PCA foreground removal performs well within the defined k-thresholds, recovering the HI power spectrum well throughout the whole $k$ range. The primary limitation is the signal loss at low $k_{\parallel}$, which is a consequence of the assumption that the foregrounds occupy lower values of $k$. This assumption breaks down when the HI signal has correlated structure on scales comparable to the foreground lengthscale and can be seen here where parts of the cosmological signal is also removed. 

This analysis would be much more difficult is a more realistic dataset were used or if a more realistic foreground model were used. The foreground signal used in this analysis assumes a single spectral index across the sky, whereas real foregrounds have spatially varying indices which would make PCA separation more difficult.

% [1] https://iopscience.iop.org/article/10.1088/2041-8205/763/1/L20
% [2] https://iopscience.iop.org/article/10.1088/0004-637X/815/1/51
In the case of more complex data with components more representative of true observations, another method which could be considered instead of applying minimum $k$ thresholds to mitigates signal loss is to use a transfer function [1], [2]. Instead of excluding the modes which cause the most signal loss, a transfer function is able to keep all $k$ modes but corrects for the signal loss at each one. Instead of completely removing low $k$ modes, it attempts to boost them by the expected suppression factor. However, while this preserves cosmological information, a limitation is that it amplifies noise and foreground residuals while the signal is being boosted. Given the significant foreground residuals present at low $k_{\parallel}$ in this analysis, applying a transfer function correction would amplify both the HI signal and the residual foreground contamination simultaneously, making the k-threshold approach the more conservative and appropriate choice here.

%Additionally, this analysis only used a small range of frequencies whereas other analysis' use a broader range. While this simplifies the foreground cleaning, this smaller range allowed for 100 realisation to be run, allowing for the foreground cleaning and cosmic variance to be taken into consideration and error bars to be quantified and plotted on the power spectra, whereas other simulation studies have a larger frequency range but only show a single realisation.

\subsection{Gaussian Process Regression}

\begin{figure*}[h!]
    \centering
    \includegraphics[width=0.71\linewidth]{outputs/week8/final_corner_plot_values.pdf}
    \caption{Posterior distribution for the MCMC sampling from GPR foreground cleaning using the exponential kernel to model the 21\,cm cosmological signal and the spectral index modified RBF kernel to model the foreground. The panels show the 1D marginals and the 2D joint constraints. The contours corresponds to the 1, 2, and 3$\sigma$ credible regions.}
    \label{fig:final_corner_plot_values}
\end{figure*}

\begin{figure*}[h!]
    \centering
    \includegraphics[width=0.71\linewidth]{outputs/week8/mcmc_chain_with_sp_indx_fg_no_noise.pdf}
    \caption{Trace plots for the MCMC sampling from GPR foreground cleaning using the exponential kernel to model the 21\,cm cosmological signal and the spectral index modified RBF kernel to model the foreground.}
    \label{fig:mcmc_chain_with_sp_indx_fg_no_noise}
\end{figure*}

Here, we will discuss the results of foreground removal using GPR. Different combinations of the kernels functions, with different initial guesses, bounds, and optimisation algorithms were tested and their LMLs compared to find the one which fit the data the best. The optimisation algorithm which found to provide the best results within the minimise function from the optimise module of SciPy was the Nelder-Mead algorithm.

The exponential function (\Cref{eq:exp}) was found to fit the 21\,cm signal the best in all cases. Additionally, the spectral index modified RBF kernel (\Cref{eq:RBF_sp_indx}) performed the best in comparison to the normal RBF kernel (\Cref{eq:RBF}), as expected as it better models the smooth foregrounds (see \Cref{fig:rbf_kernel_w_sp_indx_vs_foreground_covariance}). It was this kernel combination which gave the highest Bayes Factor and was used for MCMC Sampling.

First, we will go over the hyperparameter estimations that came as a result of the MCMC Sampling  using the kernel combination described above along with the MCMC parameter convergence. Then, we will discuss the results from the foreground removal which is the 1D and cylindrical power spectrum. Finally, we will finish with a review on the overall performance of GPR.

\subsubsection{Hyperparameter Estimation}

\begin{table*}[bt!]
    \centering
    \renewcommand{\arraystretch}{2}
    \caption{Median and $1\sigma$ estimates for the hyperparameters.}
    \label{tab:hyperparameters}
    \begin{tabular}{|l|ccccc|}
        \hline
        Hyperparameter & $\sigma_{HI}$ [mK$^2$] & $\sigma_{fg}$ [mK$^2$] & $l_{HI}$ [MHz] & ${l_{fg}}$ [MHz] & $\gamma$ \\
        \hline
         & 2.05 $^{+0.24}_{-0.22}\times10^{-4}$ & 0.21 $^{+0.14}_{-0.09}$ & 1.37 $^{+0.36}_{-0.30}$ & 780.37 $^{+156.62}_{-234.79}$ & -2.92 $^{+0.91}_{-1.01}$ \\
         \hline
         \hline
        Prior & $\mathcal{U}$ (1 $\times 10^{-6}$, 0.5)& $\mathcal{U}$ (1$\times 10^{-6}$, 0.5) & $\mathcal{U}$ (0.01, 2.0) & $\mathcal{U}$ (10.0, 1000.0) & $\mathcal{U}$ (-5.0, -1.0) \\
        \hline
    \end{tabular}
\end{table*}

\begin{figure*}[t!]
    \centering
    \includegraphics[width=1\linewidth]{outputs/week8/kernel_comparison_mcmc_with_sp_indx_no_noise.pdf}
    \caption{Comparison of the fitted GPR kernels against the true component covariances. From left to right: the fitted HI kernel (exponential), the true HI covariance, the fitted HI kernel (spectral-index-modified RBF), and the true foreground covariance.}
    \label{fig:kernel_comparison_mcmc_with_sp_indx_no_noise}
\end{figure*}

The posterior distributions of the GPR hyperparameters were explored using MCMC sampling, with the resulting corner plot shown in \Cref{fig:final_corner_plot_values}. The model includes five hyperparameters: the amplitude and the length scale of the HI signal kernel and the foreground kernel ($\sigma_{HI}$, $\sigma_{fg}$, ${l_{HI}}$, ${l_{fg}}$), and the foreground spectral index ($\gamma$). All together, these parameters control how the GPR model separates the spectrally smooth foreground from the HI signal which varies rapidly with frequency. A summary of the recovered hyperparameter values can be seen in \Cref{tab:hyperparameters}. 

The HI amplitude $\sigma_{HI}$ recovered is the most tightly constrained out of all the hyperparameters and suggests that the GPR model is correctly identifying the amplitude of the HI fluctuations even in the presence of foreground contamination. In each combination of the different kernel optimisations run before the MCMC sampler, $\sigma_{HI}$ was the most consistently recovered each time.

The recovered HI lengthscale $l_{HI}$ is moderately constrained. While this posterior peaks, it also has a tail towards larger values which is not as expected. Since this parameter controls how smooth the signal is across frequency and a smaller value results in more rapid variations, this is not consistent with the HI signal. It is interesting to see the extremely diagonal covariance plot between $\sigma_{HI}$ and $l_{HI}$, which is a contrast to all the other covariance plots, highlighting these two parameters are highly correlated.

On the other hand, the foreground amplitude $\sigma_{fg}$is considerably less well constrained. This broad and asymmetric posterior reflects a degeneracy in foreground modelling where the total foreground component can be reproduced by different combinations of amplitude and correlation length. This means that it is difficult to confine these parameters independently.

The foreground lengthscale $l_{fg}$ recovered has a particularly large and asymmetric uncertainty. The posterior rises steeply and has a long tail towards higher values, suggesting the data is consistent with a wide range of $l_{fg}$. This is an issue which kept on arising upon trying different optimisation methods. No matter what initial guesses or range of bounds were used, $l_{fg}$ would consistently hit the upper bound of the range every time, even if the bound was made smaller or wider. Physically, this large $l_{fg}$ corresponds to foreground emissions that is highly correlated in frequency and therefore very smooth which is consistent with what we expect with synchrotron-dominated foregrounds. However, the poor constraint indicates that the frequency range used may be insufficient to fully characterise $l_{fg}$. A wider bandwidth may be beneficial for this hyperparameter.

Finally, the foreground spectral index $\gamma$ is recovered with a broad posterior and the injected true value of -2.7 being within the uncertainties. While this median value is consistent with the expected behaviour, the wide posterior reflects the difficulty of constraining the spectral shape which is possibly due to the relatively narrow frequency range. The diffuse joint posterior for $l_{fg}$ and $\gamma$ in the corner plot suggests that the data cannot distinguish between the different combinations of these hyperparameters. This physically makes sense as a steeper $\gamma$ and a longer $l_{fg}$ can produce similar foreground covariance structures, resulting in these hyperparameters being difficult to constrain independently.

The MCMC trace plots for all five hyperparameters are shown in \Cref{fig:mcmc_chain_with_sp_indx_fg_no_noise}. Trace plots are important in MCMC as they verify that the chains have converged and they explore the parameter space efficiently. Each coloured line represents an individual walker's trajectory through parameter space over 10,000 steps. For a well-converged sampler where the walkers efficiently explore the posterior distribution, the traces should be dense bands that are stationary, meaning the chain has converged around a mean with variance variance with no upward or downward trends. 

For all hyperparameters, the plots demonstrate good mixing across all walkers. The mean acceptance fraction for the burn-in and sampling were both approximately 0.46 which shows there was an optimal balance between step size and efficient exploration of the posterior. $\sigma_{HI}$ is the most well well converged and stable, consistent with the constrained posteriors seen in the corner plot. However, $l_{HI}$, $\sigma_{fg}$, $l_{fg}$, and $\gamma$ show the walkers explore a wide range of values throughout the run, often hitting the lower and upper bounds throughout the run. While these chains do not show any systematic drift and appear stationary, this large spread is consistent with the broad, poorly constrained posteriors seen in the corner plot. This confirms the degeneracies the foreground model posses and that the difficulty in constraining these hyperparameters is a feature of the posterior landscape rather than the sampler failing to converge. The fitted kernels, along with the true covariance signals, can be seen in \Cref{fig:kernel_comparison_mcmc_with_sp_indx_no_noise}.

\subsubsection{Cylindrical Power Spectra}

\begin{figure*}[t!]
    \centering
    \includegraphics[width=1\linewidth]{outputs/week9/gpr_cylindrical_scipy_with_no_noise_and_sp_indx.pdf}
    \caption{The cylindrical power spectra for the GPR foreground cleaning method, with the HI signal, the residual total signal (HI signal and foreground) after foreground removal, and the ratio between residual and the HI signal, from left to right.}
    \label{fig:gpr_cylindrical_scipy_with_no_noise_and_sp_indx}
\end{figure*}

\begin{figure}
    \centering
    \includegraphics[width=1\linewidth]{outputs/week9/final_power_spectrum_1.pdf}
    \caption{The logarithmic spherically averaged 1D power spectrum from of the HI signal with no foreground removal (black), and the residual signal after GPR foreground removal (red). This was run for 1,000 realisation from the MCMC samples, where the datapoints are the median values and the 16th and 84th percentiles were used as the error bars for the residual signal.}
    \label{fig:final_power_spectrum_1}
\end{figure}

The cylindrical power spectrum provides some more information on the foreground removal by separating modes along the line of sight ($k_{\parallel}$) from transverse modes ($k_{\perp}$). \Cref{fig:gpr_cylindrical_scipy_with_no_noise_and_sp_indx} shows the true HI power spectrum, the residual signal from the GPR clean, and their ratio. 

The ratio panel is the clearest diagnostic on the success of the foreground cleaning. We can see that throughout the panel, the ratio is close to unity while at low $k_{\parallel}$ and high $k_{\perp}$, at the smaller spatial scale, there is cosmological signal loss. Additionally, there are $k$ modes at low $k_{\perp}$, suggesting under-cleaning of the large spatial scale. 

\subsubsection{1D Power Spectrum}

The spherically averaged log 1D power spectrum before and after GPR foreground removal is shown in \Cref{fig:final_power_spectrum_1}. The black curve represents the true HI power spectrum whereas the red dashed curve shows the recovered power spectrum after GPR cleaning.

At high $k$ values, the residual power spectrum lies exactly on the true HI signal, showing a perfect recovery. However, at lower $k$, which refers to large scales, values the signal deviates largely from the true signal. This is consistent with the fact that foregrounds are most dominant at low frequencies and so there is some residual foregrounds left at low $k$, indicating the foreground removal is incomplete.

It is important to note that these power spectra consist of 100 realisation from data generated from one seed. To quantify uncertainty, 1,000 parameter sets being drawn at random from the MCMC posterior chain. The median values were used as the datapoints while the 16th and 84th percentiles were used as the error bars. Hence, there are only error bars seen on the residual signal where the foreground cleaning was performed 1,000 from different hyperparameter values. 

From these error bars, we can see they are incredibly small at large $k$, but become increasingly large at low $k$. 

\subsubsection{Overall Performance}

The GPR foreground removal demonstrates partial but incomplete recovery of the HI power spectrum. This method successfully reduced the foreground contamination, and the recovered hyperparameters are broadly consistent with the injected true values or what is physically expected. However, residual foreground remains, particularly at low $k_{\perp}$ and high $k_{\parallel}$. This is likely due to the poorly constrained foreground lengthscale and the narrow frequency bandwidth of the simulated dataset. Both of these may have contributed and limited GPR's ability to clean and identify the two components individually. 
%\section{Discussion}

The performance of the PCA foreground removal method was evaluated by comparing the recovered residual signal to the HI signal. By removing $N_{FG} = 3$, the majority of the foreground, which is four ordered of magnitude larger than the 21\,cm signal, was removed. 

The ratio plot in \Cref{fig:cylindrical_power_spectrum_hi_signal} seen on the right reveals the outcome of the foreground cleaning. Due to the foregrounds being spectrally smooth, they occupy the low-$k_\parallel$ region of the power spectrum which corresponds to the large-scale frequency of the signal. However, it can be seen that low values of $k_\parallel$ have significant signal loss observed below $k_\parallel = 0.1$, indicating the 21\,cm signal cannot be separated from the foreground contribution in these modes. %Similarly, there is overlap at low $k_\perp$ due to the smoothness of the foregrounds. 

To address this, minimum thresholds of $k_\perp = 0.025$ and $k_\parallel = 0.1$ were defined and foreground cleaning of the signal was redone excluding these values to ensure minimal 21\,cm signal loss. After applying these thresholds, the resulting 1D power spectrum can be seen in \Cref{fig:log_power_spectrum_hi_signal}. The resulting residual signal shows strong agreement with the 21\,cm signal at the high $k$. This suggests that PCA is effective at recovering the high-frequency fluctuations of the 21\,cm signal and is physically expected since high $k_\parallel$ correspond to small-scale frequency structure that the foregrounds do not contaminate. However, the residual signal is slightly above the expected HI signal at lower $k_\parallel$, suggesting there are residual foregrounds that PCA did not fully capture and the cleaning is incomplete.

%Therefore, while PCA is an effective blind foreground removal method, its main limitation is the removal of entire modes which results in the observed signal loss at low $k_\parallel$. 

%Therefore, while PCA is an effective blind foreground removal method for recovering high-frequency 21\,cm fluctuations, its main limitation is the loss of entire modes at low k∥ and incomplete foreground separation near the foreground-dominated boundary.
\section{Conclusion}

This section will have the conclusion but is not written yet.

%------------------------------------------------

%----------------------------------------------------------------------------------------
%	 REFERENCES
%----------------------------------------------------------------------------------------
% Apparently we should use %\autocite{} instead of \cite{} it makes citations format better.

\renewcommand{\refname}{Bibliography}
\addcontentsline{toc}{section}{Bibliography}
% \bibliographystyle{plain}
%\bibliographystyle{unsrt}
\printbibliography

\onecolumn
\appendix
\addcontentsline{toc}{section}{Appendix}
%\section{GPR results without spectral index}

\begin{table*}[htb!]
    \centering
    \renewcommand{\arraystretch}{2}
    \caption{Median and $1\sigma$ estimates for the hyperparameters.}
    \label{tab:hyperparameters}
    \begin{tabular}{|l|cccc|}
        \hline
        Hyperparameter & $\sigma_{HI}^2$ [mK$^2$] & $\sigma_{fg}^2$ [mK$^2$] & $\mathcal{l_{HI}^2}$ [MHz] & $\mathcal{l_{fg}^2}$ [MHz] \\
        \hline
         & 0.0013 $^{+0.0001}_{-0.0001}$ & 0.3084 $^{+0.1232}_{-0.1137}$ & 0.0752 $^{+0.0292}_{-0.0381}$ & 549.4692 $^{+255.4766}_{-216.1621}$ \\
         \hline
         \hline
        Prior & $\mathcal{U}$ (0.000001, 0.5)& $\mathcal{U}$ (0.000001, 0.5) & $\mathcal{U}$ (0.01, 2.0) & $\mathcal{U}$ (10.0, 1000.0) \\
        \hline
    \end{tabular}
\end{table*}

\begin{figure*}[hbt!]
    \centering
    \includegraphics[width=1\linewidth]{outputs/week8/cylindrical_power_spectrum_mcmc_without_sp_indx_fg_2.pdf}
    \caption{Enter Caption}
    \label{fig:placeholder}
\end{figure*}

\begin{figure*}[hbt!]
    \centering
    \includegraphics[width=0.7\linewidth]{outputs/week8/mcmc_corner_plot_without_sp_indx_fg_2.pdf}
    \caption{Enter Caption}
    \label{fig:placeholder}
\end{figure*}

\begin{figure*}[hbt!]
    \centering
    \includegraphics[width=0.7\linewidth]{outputs/week8/mcmc_chain_without_sp_indx_fg_2.pdf}
    \caption{Enter Caption}
    \label{fig:placeholder}
\end{figure*}
\input{documents/lab_report/appendices/appendix_b}
\newpage
%----------------------------------------------------------------------------------------

\end{document}
